% ---------------------------------------------------------------------
% Section 5: Applications and examples.
% Sources: LRMQG.tex table of POMDP examples (inserted as-is per the
% authors' instruction, including the authors' margin notes); Normal/
% Poisson/Uniform garbling constructions; two-state technology adoption
% example with primitives recovered from 2statesTechAdoptionClaude.xlsx.
% ---------------------------------------------------------------------
\section{Applications}\label{sec:applications}

\subsection{Three Monotone POMDPs}\label{sec:three-models}

Table~\ref{tab:pomdp-examples} illustrates the scope of the monotone
POMDP framework with three models from the literature: the technology
adoption model of \citet{ulu2009uncertainty}, the machine replacement
model of \citet{white1979optimal} (see also \citealt{lovejoy1987some}),
and the prostate biopsy referral model of \citet{zhang2012optimization}.
In each model, the state, actions, and signals are ordered, the value
function is monotone in the beliefs, and the optimal policy is
monotone: as beliefs about the technology improve, the DM moves toward
adoption; as beliefs about the machine's condition worsen, the DM moves
toward replacement; and as beliefs point toward cancer, the DM moves
toward biopsy. Theorem~\ref{thm:main} applies to each model: improving
the signal process in the LR-better sense increases the optimal value
function. \Claude{Transitional text drafted; the table is inserted
as-is from your notes, including the two ``?'' cells and your margin
notes, per your instruction.}

\begin{landscape}
\begin{table}[h!]
\footnotesize
\centering
\caption{POMDP examples: \citet{ulu2009uncertainty},
\citet{white1979optimal}, \citet{zhang2012optimization}.}
\label{tab:pomdp-examples}
\renewcommand{\arraystretch}{1.4}
\begin{tabular}{>{\bfseries}p{2.5cm} p{6.5cm} p{6.5cm} p{5.5cm}}
\toprule
 & \textbf{Technology Adoption} & \textbf{Machine Replacement} & \textbf{Prostate Biopsy} \\
State of the world
    & $\theta_k \in \Theta \subseteq \Re, $ \newline technology benefit in period $k$
    & $s_t \in S=\{0,\ldots,n=\text{good as new}\}$, \newline machine state in period $t$
    & $s_t\in \{NC, C, T, M, D\}$, $NC$: no cancer, $C$: undetected prostate cancer,
    $T$: Treated, nonmetastatic (obs.), \newline $M$: metastasis (obs.), $D$: Death (obs.)\\
\midrule
Prior
  &$\pi(\theta_k)$, \newline continuous probability distn.
  & $\pi(s_t)$, \newline discrete probability distn.
  & $\pi_t \in [0,1]$, \newline probability of undetected cancer \\
\midrule
Signal process
    & $L(x|\theta_k)$ probability of observing signal $x$ given technology state $\theta_k$ \newline
    satisfies the MLR property
    &$r_{jk}^a$: probability of observing $k\in O$ given state $s_{t+1}=j$ and action $a_t=a$
  & $q_t(l_t|s_t)$: probability of observing PSA result $l_t \in \{1,\ldots,m\}$
    given state $s_t$ \\
\midrule
Actions $(a_k)$
  & $a_k = $ technology choice in period\newline
    $\in \{$adopt, investigate, reject$\}$\newline
    where adopt $\succ$ investigate $\succ$ reject
  & $a_t \in \{a:$ produce, $a^\prime:$ replace$\}$\newline
  where  $a^\prime\succ a$
  & $a_t = $ biopsy decision in period  $\in \{B, W\}$\newline
   where $W\succ B$\\
\midrule
Rewards\newline$(r_k(\theta_k, a_k))$
  & $r_k(\theta_k, a_k)$\newline
    $= \begin{cases} A\theta_k - K & \text{if } a_k = \text{adopt} \\ -c & \text{if } a_k = \text{investigate} \\ 0 & \text{if } a_k = \text{reject} \end{cases}$
  & $g(i,a)$: reward from action $a$ in state $s_i$, \newline
    has isotone differences
  & $\bar{r}_t(s_t, a_t)$\newline
    $= \begin{cases} 1 & \text{if } a_t = W \\ 1-\mu & \text{if } a_t = B, s_t=NC \\ 1-\mu-f_\epsilon & \text{if } a_t = B, s_t=C\end{cases}$ \newline{\CU{There are more states.}}\\
\midrule
Expected Rewards\newline$(r_k(\pi, a_k))$
  & $r_k(\pi, a_k)=\int_{\theta \in \Theta} r_k(\theta,a_k) \pi(\theta) \dif \theta$
  & $\sum_{i\in S} \pi_i g(i,a)$
  & $r_t(\pi_t, a_t)$\newline
    $= \begin{cases} 1 & \text{if } a_t = W \\ R_t(\pi_t) & \text{if } a_t = B\end{cases}$\\
\midrule
Transitions\newline$(\Pi(\theta_k;\pi, x,a_k))$
  & $\Pi(\theta_k;\pi, x,a_k)$\newline
    $= \begin{cases} ? & \text{if } a_k = \text{adopt or reject} \\ \frac{ L(x|\theta,a) \eta(\theta; \pi, a)}{f(x;\pi, a)}& \text{otherwise} \end{cases}$\newline
    where $\eta(\theta_k; \pi, a)=\int_{\theta_k \in \Theta} g(\theta_{k-1}|\theta_k, a) \pi(\theta_k) \dif \theta_k$ and $f(x;\pi, a)= \int_\theta L(x|\theta,a) \eta(\theta; \pi, a)\,\dif \theta$
  & $T(\pi, a, k)$\newline
   $=\begin{cases} ? & \text{if } a_k = \text{produce} \\ e_n \footnotesize{\text{(a unit vector with a 1 in the $n^{th}$ position)}}& \text{if } a_k=\text{replace} \end{cases}$
  & $\pi_{t+1}(s_{t+1})=\frac{q_{t+1}(l_{t+1}|s_{t+1})\sum_{s_t \in S}p_t(s_{t+1}|s_t, a_t)\pi_t(s_t)}{\sum_{s_{t+1}\in S}q_{t+1}(l_{t+1}|s_{t+1})\sum_{s_t \in S}p_t(s_{t+1}|s_t, a_t)\pi_t(s_t)}$\newline \CU{we can use something like our notation here.}\\
\midrule
Properties of the value function
  & $V_k^*(\pi)$ is $LR$-increasing and convex in $\pi$, \newline
    $V_k(\pi, a_k)$ satisfies $LR$ increasing differences.
  & For each $s_k$, $v(s_k, p_k; \sigma)$ is increasing\newline
    and convex in $p_k$ and increasing in\newline
    $\sigma$
  & $v_t(\pi_t)$ is decreasing and convex in $\pi_t$,\newline $R_t(\pi_t)$ linear and decreasing in $\pi_t$, \newline Control-limit type policy optimal \\
\bottomrule
\end{tabular}
\end{table}
\end{landscape}

\subsection{Constructing LR-Better Signal Processes}\label{sec:examples}

We next present three families of examples in which an LR-better
garbling arises naturally: a Normal signal with a state-dependent
bias, a Poisson signal with state-dependent thinning, and a Uniform
signal with state-dependent scaling. In each case, the garbled process
would typically \emph{not} be comparable to the original in
Blackwell's order because the added noise depends on the state, yet
the LR-better order applies and, by Theorem~\ref{thm:main}, a monotone
POMDP has a higher value function with the original signal process.

\subsubsection{Normal Likelihood with a State-Dependent Bias.}
\label{sec:normal-example}
Let $F(x|\mu, \sigma^2)$ be a Normal signal-generating process with
mean $\mu$ and variance $\sigma^2$, where the state of the world is
the unknown mean $\mu$. This likelihood satisfies the MLR property.
Suppose that we add a bias term that depends on $\mu$ and observe
\begin{equation}
Y = X + d(\mu) + Z \label{eq:normal-garbling}
\end{equation}
instead, where $d(\mu)$ is a decreasing function of $\mu$ and $Z$ is a
Normally distributed error term with mean $0$ and variance $\nu^2$,
independent of $X$. The garbling kernel is then
\begin{equation}
H(y|x, \mu) \sim N(x+d(\mu),\, \nu^2) \ ,
\end{equation}
a conditional probability distribution that is LR-decreasing in $\mu$
because $d(\mu)$ is decreasing; for example, $d(\mu) =
-\delta\mu/\sigma$ for some $\delta > 0$. Given the definition of the
biased signal process $\iY$, we have
\begin{equation}
G(y|\mu) \sim N(\mu+d(\mu),\, \sigma^2+\nu^2) \ ,
\end{equation}
which also satisfies the MLR property provided that $\mu + d(\mu)$ is
increasing; with $d(\mu) = -\delta\mu/\sigma$, this holds if
$\delta \leq \sigma$. Then $\iX \LRI \iY$: since both likelihoods
satisfy the MLR property, Theorem~\ref{thm:main} implies that a
monotone POMDP has a higher value function with likelihood $F$ than
with $G$. If $d(\mu)$ is constant, the noise is state independent and
$F$ Blackwell-dominates $G$; a strictly decreasing $d$ takes the
comparison outside Blackwell's order but keeps it within the LR-better
order.

The bias term has a natural interpretation: signals about better
technologies may be systematically dampened---for example, reviews of
genuinely superior products may understate their advantage---and such
skepticism-type distortions degrade information in exactly the
reversely monotone way the LR-better order captures.
\Claude{Interpretation sentence drafted by me---replace or delete if
it does not match the story you want to tell.}

\subsubsection{Poisson Likelihood with State-Dependent Thinning.}
\label{sec:poisson-example}
Let $F(x|\lambda)$ be a Poisson distribution with rate $\lambda$,
which has the MLR property in $\lambda$. Suppose that, instead of
observing the count $X$, each arrival is observed independently with
probability $d(\lambda)$. The garbling distribution $H(y|x,\lambda)$
is then Binomial with $x$ trials and success probability
$d(\lambda)$; it is LR-decreasing in $\lambda$ if $d(\lambda)$ is
decreasing. We now show that, after this thinning, the resulting
signal distribution is Poisson with rate $d(\lambda)\lambda$:
\begin{align}
g(y|\lambda) &= \sum_{x=y}^\infty h(y|x, \lambda)\, f(x|\lambda) \\
&= \sum_{x=y}^\infty \binom{x}{y}\, [d(\lambda)]^y
  [1-d(\lambda)]^{x-y}\, \frac{e^{-\lambda}\lambda^x}{x!} \\
&= \frac{e^{-\lambda}[d(\lambda)]^y}{y!} \sum_{x=y}^\infty
  \frac{\lambda^x [1-d(\lambda)]^{x-y}}{(x-y)!} \\
&= \frac{e^{-\lambda}[d(\lambda)]^y}{y!} \sum_{z=0}^\infty
  \frac{\lambda^{z+y} [1-d(\lambda)]^{z}}{z!} \\
&= \frac{e^{-\lambda}[d(\lambda)\lambda]^y}{y!} \sum_{z=0}^\infty
  \frac{\left(\lambda[1-d(\lambda)]\right)^{z}}{z!} \\
&= \frac{e^{-\lambda}[d(\lambda)\lambda]^y\,
  e^{\lambda(1-d(\lambda))}}{y!} \\
&= \frac{e^{-\lambda d(\lambda)}[d(\lambda)\lambda]^y}{y!} \ .
\end{align}
The second equality uses the definitions of the Binomial and Poisson
distributions; the third uses $\binom{x}{y} = \frac{x!}{(x-y)!\,y!}$;
the fourth substitutes $z = x-y$; the fifth collects the powers of
$\lambda$; and the sixth uses $\sum_{z=0}^\infty x^z/z! = e^x$. The
resulting Poisson distribution with rate $d(\lambda)\lambda$ has the
MLR property if $d(\lambda)\lambda$ is increasing in $\lambda$. Thus,
if $d$ is decreasing but $d(\lambda)\lambda$ is increasing, we have
$\iX \LRI \iY$ with both processes satisfying the MLR property, and
Theorem~\ref{thm:main} applies. This is a model of state-dependent
undercounting: arrivals from better processes (higher $\lambda$) are
individually less likely to be recorded.

\subsubsection{Uniform Likelihood with State-Dependent Scaling.}
\label{sec:uniform-example}
Let $F(x|\theta)$ be a Uniform distribution on $[0, \theta]$, which
has the MLR property in $\theta$. Suppose that instead of observing
$X$ we observe
\begin{equation}
Y = d(\theta) X \ .
\end{equation}
Then $Y|\theta$ is Uniform on $[0, d(\theta)\theta]$, which has the
MLR property if $d(\theta)\theta$ is increasing in $\theta$. For this
transformation, the garbling distribution is deterministic:
\begin{equation}
H(y|x,\theta) = \left\{
\begin{array}{cl}
1 & \mbox{if } y = x\, d(\theta)\ , \\
0 & \mbox{otherwise}\ .
\end{array}\right.
\end{equation}
To see that $H(y|x,\theta)$ is LR-decreasing when $d(\theta)$ is
decreasing in $\theta$, consider the inequality
\begin{equation}
H(y^2|x,\theta^2)\, H(y^1|x,\theta^1) \leq
H(y^2|x,\theta^1)\, H(y^1|x,\theta^2) \label{eq:uniform-LR}
\end{equation}
for $y^1 \leq y^2$ and $\theta^1 \leq \theta^2$. If
$y^1 = x\, d(\theta^2) \leq y^2 = x\, d(\theta^1)$, the inequality
holds because the left side is $0$ and the right side is $1$; for any
other selection of $y^1 \leq y^2$, both sides equal zero.
\CU{I think the above example would work for $Y = X + d(\theta)$. I
tried more general garbling distributions $H(y|x,\theta)$ that are not
deterministic. For example, we can observe $x$ or $x\,d(\theta)$ with
equal probability (does not satisfy the LR-decreasing property), or
$H(y|x,\theta)$ uniform between $x$ and $x+d(\theta)$ (is
LR-decreasing if $d(\theta)$ is decreasing, but the marginal for
$y|\theta$ is a trapezoidal distribution and it is not clear if it can
be made to have the MLR property).}

\subsection{A Two-State Technology Adoption Example: When
State-Increasing Garbling Yields No Ranking}\label{sec:tech-example}

The reversely monotone structure of the LR-better order---noise that
degrades the signal more in \emph{higher} states---is not a technical
artifact: when the state dependence of the noise runs in the opposite
direction, the two signal processes may genuinely be incomparable. We
illustrate this with a simple two-state technology adoption example.

The technology's benefit is $\theta \in \{0.1, 0.7\}$; we refer to
these as the low and high states. The DM's prior probability on the
high state is $0.4$. The DM either adopts the technology, earning
$A\theta - K$ with $A = 1$ and adoption cost $K = 0.3$ (i.e., $-0.2$
in the low state and $0.4$ in the high state), or rejects it, earning
$0$. Before deciding, the DM observes the outcome of two independent
trials of the technology, each of which succeeds with probability
$\theta$.

Under signal process $\iX$ (experiment $F$), the DM observes the
number of successes $x \in \{0,1,2\}$, so $L_\iX(x|\theta)$ is
Binomial with two trials and success probability $\theta$; this
likelihood satisfies the MLR property. Under signal process $\iY$
(experiment $G$), each success is \emph{recorded} only with
probability $d(\theta)$, with $d(0.1) = 0.4$ and $d(0.7) = 0.7$: the
probability that a success is observed \emph{increases} with the
state. The DM observes the number of recorded successes
$y \in \{0,1,2\}$, so $L_\iY(y|\theta)$ is Binomial with two trials
and success probability $\theta\, d(\theta)$ ($0.04$ in the low state
and $0.49$ in the high state), which also satisfies the MLR property.
The garbling kernel taking $x$ to $y$ is Binomial with $x$ trials and
success probability $d(\theta)$---exactly as in the Poisson thinning
example of Section~\ref{sec:poisson-example}, except that here $d$ is
\emph{increasing} in the state, so the kernel is LR-\emph{increasing}
in $\theta$ and the LR-better condition (Definition~\ref{def:LRI}(ii))
fails; indeed, even the weaker FOSD condition of monotone
quasi-garbling fails.

Figure~\ref{fig:tech-trees} shows the decision trees for the two
experiments. At the prior of $0.4$, experiment $F$ is more valuable:
the DM's expected value is $0.123$ under $F$ versus $0.109$ under $G$.
But this ranking does not hold for all priors: for sufficiently
pessimistic priors, $G$ is strictly more valuable (for example, at a
prior of $0.01$ on the high state, the DM's value is $0.0006$ under
$G$ and $0$ under $F$---under $F$, no signal is favorable enough to
justify adoption, while under $G$ the strongest signal is). The two
value functions cross as a function of the prior, so neither
experiment dominates the other. \Claude{Values computed from your
spreadsheet (2statesTechAdoptionClaude.xlsx): the value-by-prior table
shows $G$ strictly better for priors around $0.01$ and $F$ strictly
better for moderate priors, e.g., $0.1228$ vs.\ $0.10898$ at prior
$0.4$.}

Why does no ranking exist? When the observation probability increases
with the state, the signal under the signal process $G$ carries two
pieces of evidence rather than one: a recorded success indicates not
only that the trial succeeded, but also that the state was one in
which successes are more likely to be recorded. These two pieces of
evidence point in the same direction and reinforce each other, so a
favorable signal under $G$ is stronger evidence for the high state
than the corresponding signal under $F$. The price is paid at the
other end of the signal scale: silence becomes ambiguous, since it may
reflect either a failed trial or an unobserved success, and therefore
constitutes weaker evidence for the low state than it does under $F$.
The effect on beliefs is a one-sided stretch: posteriors move farther
than under $F$ after good news, but not as far after bad news. $G$ is
thus genuinely more informative for decisions that hinge on confirming
the high state and genuinely less informative for decisions that hinge
on detecting the low state, and neither experiment can dominate the
other: which one is more valuable depends on the prior beliefs.

\begin{figure}[t]
\FIGURE{%
\begin{tabular}{c@{\qquad}c}
\resizebox{0.45\textwidth}{!}{%
\begin{tikzpicture}[x=1.15cm,y=1.15cm]
\node[chance] (root) at (0,-3.6125) {};
\node[nval] at (root.south) {$0.123$};
\payoffheader{0.75}
\sigblock{root}{0}{.522\ $x=0$}{.931}{.069}{$0$}{$-0.159$}{plain}{opt}
\sigblock{root}{-2.55}{.276\ $x=1$}{.391}{.609}{$0.165$}{$0.165$}{opt}{plain}
\sigblock{root}{-5.1}{.202\ $x=2$}{.030}{.970}{$0.382$}{$0.382$}{opt}{plain}
\end{tikzpicture}}
&
\resizebox{0.45\textwidth}{!}{%
\begin{tikzpicture}[x=1.15cm,y=1.15cm]
\node[chance] (root) at (0,-3.6125) {};
\node[nval] at (root.south) {$0.109$};
\payoffheader{0.75}
\sigblock{root}{0}{.657\ $y=0$}{.842}{.158}{$0$}{$-0.105$}{plain}{opt}
\sigblock{root}{-2.55}{.246\ $y=1$}{.187}{.813}{$0.288$}{$0.288$}{opt}{plain}
\sigblock{root}{-5.1}{.097\ $y=2$}{.010}{.990}{$0.394$}{$0.394$}{opt}{plain}
\end{tikzpicture}}
\\[4pt]
(a) Experiment $F$. & (b) Experiment $G$.
\end{tabular}}%
{Decision trees for the two-state technology adoption example
($\theta\in\{0.1,0.7\}$, prior on the high state $0.4$, adoption cost
$0.3$).\label{fig:tech-trees}}%
{Squares are decision nodes, circles are chance nodes, and bold
branches mark optimal decisions; node values appear below each node.}
\end{figure}
